\section{\textbf{結果與討論}}
實驗一： 理想上兩電極中央應為零電位，但實測發現向左偏移（6.7 cm vs 7.7 cm）。經量測電阻，左半部總阻值（$55.08\,\text{M}\Omega$）比右半部（$54.54\,\text{M}\Omega$）大，導致左側分壓較多、電位降較快，這解釋了零電位線左移的原因。

實驗二： U 形電極雖然是導體，但因有開口，電力線仍會進入其中。此外，電力線集中在尖角處，印證了尖端效應，說明尖銳處的電場強度較大。

實驗三： 半圓形電極表面積較大，電荷密度 $\sigma$ 較低，使電力線分布稀疏；而左側小面積電極的電力線明顯密集。這說明了電極形狀會直接影響電場強度的分布。